\subsection{A full level-seven basis over integer residue rings}
\label{sec:gl2-minus7-generic}

\begin{proposition}[K17-MINUS7-GL2-GENERIC]
\label{prop:minus7-gl2-generic}
Let \(e\ge2\), put \(m=2e-2\) and \(d=2m+1=4e-3\), and let
\(T_2\) act on the entire space \(S_{2e+1}(\Gamma_1(7))\).
This space has a full integral Fourier basis of size \(d\), over
\(\mathbf Z[1/7]\), whose first \(d\) positive Fourier coefficients
form a lower triangular matrix with diagonal entries one.
Writing \(t_e=1+2^{2e}\), define
\[
 J_e=\det(t_e^3I-T_2^3)\in\mathbf Z[1/7].
 \tag{G7G.1}
\]
For any integer \(M>1\) with \(\gcd(M,42)=1\), this determinant
can be computed over \(\mathbf Z/M\mathbf Z\) using the full
basis and only explicitly verified unit divisions.
If its residue is a unit, then, at every prime \(\ell\mid M\),
\(t_e-\zeta a_2(f)\) is a unit at every coefficient place over
\(\ell\), for every normalized eigenform
\(f\in S_{2e+1}(\Gamma_1(7))\) and every \(\zeta\in\mu_3\).
No primality assumption on \(M\) is needed.
\end{proposition}

\begin{proof}
Use the integral weight-one generators \(z_1,z_2,z_3\) of
\cite[Lemma~2.3 and Proposition~2.4, pp.~9--11]{K17MeierOzornova2020}.
The complete graded ring is
\[
 \mathbf Z[1/7][z_1,z_2,z_3]/
       (z_1z_2+z_2z_3+z_3z_1).
 \tag{G7G.2}
\]
Set
\[
 A=z_3+z_1,\qquad B=z_1,\qquad C=z_1+z_2.
\]
This is an integral unimodular change of variables, with inverse
\(z_1=B,\ z_2=C-B,\ z_3=A-B\); its relation is \(AC=B^2\).
With \(Q=\exp(2\pi i\tau)\), the explicit series of the cited lemma give
\[
 A=1+O(Q),\qquad B=Q+O(Q^2),\qquad C=Q^2+O(Q^3).
 \tag{G7G.3}
\]
Division by the monic relation with leading monomial \(AC\)
shows that the following are a basis of the entire weight-\(m\)
piece:
\[
 A^{m-j}B^j\quad(0\le j\le m),\qquad
 B^{m-j}C^j\quad(1\le j\le m).
 \tag{G7G.4}
\]
Their leading orders are successively \(0,1,\ldots,2m\), and
their leading coefficients are one. In particular this is a
full basis over \(\mathbf Z[1/7]\), not just an independent family.
Its triangular Fourier matrix also proves directly that every
characteristic-zero form with integral Fourier coefficients at a
prime different from \(7\) has integral coordinates in this basis.

Let \(D=\eta(\tau)^3\eta(7\tau)^3\).
As established in Proposition~\ref{prop:minus7-gl2}, the eta
transformation theorem
\cite[Theorem~1.1, Remark~1.2 and Proposition~3.1]{K17SavittNewman}
gives weight three and character \(\chi_{-7}\), and \(D\) has
a simple zero at each of the six cusps of \(\Gamma_1(7)\)
and no zero in the upper half-plane. Therefore
\[
 S_{2e+1}(\Gamma_1(7))=D\,M_{2e-2}(\Gamma_1(7)).
 \tag{G7G.5}
\]
This also identifies the full integral Fourier lattices:
\(D/Q\) and its inverse belong to \(1+Q\mathbf Z[[Q]]\).
Thus \(D\) times (G7G.4) is the asserted full basis, with
successive leading orders \(1,\ldots,d\).

The diamond operator \(\delta=\langle2\rangle\) cyclically
permutes \(z_1,z_2,z_3\), by
\cite[Lemma~2.6, p.~11]{K17MeierOzornova2020}.
It fixes \(D\), since \(\chi_{-7}(2)=1\), and hence
\[
 \delta A=C,\qquad
 \delta B=C-B,\qquad
 \delta C=A+C-2B.
 \tag{G7G.6}
\]
For every form \(F\) in the full space, its Hecke series is
\[
 [Q^n]T_2F=[Q^{2n}]F+
  2^{2e}\begin{cases}[Q^{n/2}](\delta F),&2\mid n,\\0,&2\nmid n.\end{cases}
 \tag{G7G.7}
\]
This follows on every diamond eigenspace from the usual
nebentypus Hecke formula and hence on their direct sum in
characteristic zero. Formula (G7G.7) preserves integral Fourier
coefficients, and the triangular basis therefore proves lattice
stability. Its first \(d\) coefficients determine the exact
integral coordinates of \(T_2F\), by forward substitution.
The same calculation remains valid after reduction modulo any
integer \(M\) coprime to \(7\).

Over a coefficient extension containing \(\mu_3\), the commuting
matrix factorization gives
\[
 t_e^3I-T_2^3=\prod_{\zeta\in\mu_3}(t_eI-\zeta T_2).
 \tag{G7G.8}
\]
The eigenvalues of the integral operator are algebraic integers
at all primes different from \(7\). If \(J_e\) is a unit modulo
\(M\), it remains a unit at every coefficient place above
each \(\ell\mid M\). Each individual factor in (G7G.8), and
each eigenvalue difference in that factor, is integral there.
Their product is a unit, so every one is a unit. This proves
the final assertion.
\end{proof}

\begin{proposition}[K17-MINUS7-GL2-HIGHER3]
\label{prop:minus7-gl2-higher3}
Use the five residual integers \(M_{e,i}\) of
Proposition~\ref{prop:hecke-higher3}.
For the following pairs \((e,i)\), the full determinant (G7G.1)
has the displayed residue \(j\) modulo \(M_{e,i}\), and
\(\gcd(j,M_{e,i})=1\):
\[
\begin{array}{c|c|r}
e&i&j\\ \hline
20&1&4822435123\\
20&2&205886586892358540462732749\\
22&1&79920658975217\\
22&2&729501190808707622310038510098001514556\\
24&1&1165633710051053709364161260919098058495500
\end{array}
\tag{G7G.9}
\]
The respective full dimensions are \(77,85,93\).
\end{proposition}

\begin{proof}
The standalone certificate
\path{hecke_minus7/higher3/gl2_twist_residual_certificate.py}
reconstructs the three integral series from their divisor sums.
Their coefficient tables, indexed by divisors modulo \(7\),
are respectively
\[
\begin{split}
 &(0,1,-1,-2,2,1,-1),\\
 &(0,-1,2,-1,1,-2,1),\\
 &(0,2,1,1,-1,-1,-2);
\end{split}
\]
only \(z_3\) has nonzero constant coefficient, equal to one.
It reconstructs
\(D=Q\prod_{n\ge1}(1-Q^n)^3(1-Q^{7n})^3\),
checks \(AC=B^2\) to its entire precision, and constructs every
one of the \(77,85,93\) basis forms and their diamond images.
It checks that the first \(d\) positive coefficients are
unit lower triangular, and uses (G7G.7) to construct the full
matrix by forward substitution. The total precisions are
\(170,186,202\), respectively, and every coefficient through
\(d+8=85,93,101\) is independently checked after reconstruction.
The identification of the operator follows from the complete
space and the triangular matrix; the additional coefficients
are consistency checks.

The program computes \(T_2^3\) by matrix multiplication over
\(\mathbf Z/M\mathbf Z\) and evaluates (G7G.1) by elimination.
Every division uses a pivot whose gcd with \(M\) is checked
to be one; its inverse is obtained by the extended Euclidean
algorithm and verified by multiplication. A missing unit pivot
would be reported as unresolved, with the current column-ideal
gcd, rather than as a zero determinant. In the five calculations
all pivots are units. The recorded residues are exactly
(G7G.9), and direct integer gcd computations give one in
every case. The accompanying JSON preserves the entire
triangular matrices, initial Hecke coefficients, full Hecke
matrices, every determinant pivot, and all five residues.
This is a computation in integer quotient rings throughout;
neither primality of \(M\) nor the existence of a chosen
root of unity modulo \(M\) enters it.
\end{proof}
